\section*{الفصل \ref{ch:g12:climate-change} --- تغير المناخ: الأرصاد والأسباب والإسقاطات}

\begin{solution}{exo:g12:climate-change:1}
التغير الذي يحدثه في صافي تدفق الطاقة الداخل إلى النظام المناخي
عند قمة الغلاف الجوي، قبل أن تستجيب درجة حرارة السطح؛ ويُقاس
بوحدة \unit{W/m^2}. الموجب: يكسب الكوكب طاقة ويسخن؛
والسالب: يفقد طاقة ويبرد.
\end{solution}

\begin{solution}{exo:g12:climate-change:2}
أعلى ما يكون في مايو، وأدنى ما يكون في سبتمبر. ففي الربيع والصيف يأخذ الغطاء النباتي
في نصف الكرة الشمالي، حيث تنمو معظم نباتات اليابسة، من ثاني أكسيد الكربون
بالتمثيل الضوئي أكثر مما تطلقه النباتات والتربة؛ وفي الخريف
والشتاء يغلب التحلل والتنفس فيرتفع التركيز من جديد.
\end{solution}

\begin{solution}{exo:g12:climate-change:3}
\emph{مرجح جدًا}: احتمال من 90 إلى 100\,\%؛ و\emph{مرجح}: من 66 إلى
100\,\%. و$[0.95 \text{, } 1.20]$ مدى مرجح جدًا: أي فرصة 90\,\% لأن
تقع القيمة الحقيقية بين 0.95 و1.20.
\end{solution}

\begin{solution}{exo:g12:climate-change:4}
91\,\%. فالماء يتمدد حين يسخن: والمحيط الأدفأ يشغل حجمًا أكبر، فيرتفع
سطحه حتى لو لم يُضف إليه ماء (ويضيف ذوبان الجليد القاري ماءً فوق
ذلك).
\end{solution}

\begin{solution}{exo:g12:climate-change:5}
بخار الماء: يحمل الهواء الأدفأ مزيدًا من بخار الماء، وهو غاز دفيئة، فيزيد
الاحترار. الثلج والجليد: الذوبان يكشف سطوحًا أدكن، تمتص
مزيدًا من ضوء الشمس. والاستجابة السالبة: الكوكب الأدفأ يشع مزيدًا من الأشعة تحت الحمراء
إلى الفضاء، فيعود إلى التوازن.
\end{solution}

\begin{solution}{exo:g12:climate-change:6}
$427.35 - 315.98 = \qty{111.37}{ppm}$؛ و$111.37/315.98 \approx 0.352$، أي
زيادة بنحو 35\,\%. وعلى مدى 66 سنة: $111.37/66 \approx
\qty{1.7}{ppm}$ في السنة في المتوسط.
\end{solution}

\begin{solution}{exo:g12:climate-change:7}
$F = 3.93 \times \ln(427.35/278)/\ln 2 = 3.93 \times 0.4300/0.6931 \approx
\qty{2.44}{W/m^2}$.
\end{solution}

\begin{solution}{exo:g12:climate-change:8}
$3.7 \times 30 = \qty{111}{mm}$، أي نحو \qty{11}{cm}. و$3.7/1.3 \approx 2.8$:
أسرع بنحو ثلاث مرات تقريبًا.
\end{solution}

\begin{solution}{exo:g12:climate-change:9}
$\Delta T = 2.72/1.16 \approx \qty{2.3}{\degreeCelsius}$. فالمناخ ليس
في حالة توازن: فالمحيط، الذي يأخذ نحو 90\,\% من الطاقة الإضافية،
يسخن ببطء، ولذلك لم يلحق السطح بعدُ \omterm{def:g12:climate-change:forcing}{بالتأثير الإشعاعي} (فلا يزال
النظام يكسب نحو \qty{0.8}{W/m^2}).
\end{solution}

\begin{solution}{exo:g12:climate-change:10}
لا تفسر الشمس والبراكين إلا \qty{0.1}{\degreeCelsius} على الأكثر من
احترار قدره \qty{1.06}{\degreeCelsius}: فلا يمكن أن تكون سببه. والسبب الطبيعي
الآخر هو التقلبية الداخلية للمناخ (تقلباته
الطبيعية، كتبادل الحرارة بين المحيط والغلاف الجوي)؛ ويقدّر
التقييم مساهمتها بمقدار $-0.2$ إلى $+\qty{0.2}{\degreeCelsius}$،
وهي أيضًا أصغر بكثير من أن تفسره.
\end{solution}

\begin{solution}{exo:g12:climate-change:11}
دون \qty{2}{\degreeCelsius} على نحو مرجح جدًا: SSP1-1.9 وحده (مداه المرجح
جدًا من 1.0 إلى 1.8). وفوقها على نحو مرجح جدًا: SSP2-4.5 (من 2.1 إلى 3.5)، وSSP3-7.0 (من 2.8 إلى 4.6)
وSSP5-8.5 (من 3.3 إلى 5.7). أما SSP1-2.6 (من 1.3 إلى 2.4) فلا يكون فيه أي منهما مرجحًا جدًا.
\end{solution}

\begin{solution}{exo:g12:climate-change:12}
\omterm{def:g12:climate-change:forcing}{تأثير} الهباء الجوي $-0.22 - 0.84 = \qty{-1.06}{W/m^2}$؛ ومن دونه يصير
المجموع $2.72 + 1.06 = \qty{3.78}{W/m^2}$. ويزداد الاحترار عند التوازن
بمقدار $1.06/1.16 \approx \qty{0.9}{\degreeCelsius}$.
\end{solution}

\begin{solution}{exo:g12:climate-change:13}
لأن $\ln(2C) - \ln C = \ln 2$ أيًّا كانت $C$،
\[
  F(556) - F(278) = \frac{3.93\,\ln 2}{\ln 2} = \qty{3.93}{W/m^2}, \qquad
  F(1112) - F(556) = \frac{3.93\,\ln 2}{\ln 2} = \qty{3.93}{W/m^2}:
\]
أي \omterm{def:g12:climate-change:forcing}{التأثير} نفسه لكل مضاعفة. والمجموع عند \qty{1112}{ppm}: $2 \times 3.93 =
\qty{7.86}{W/m^2}$.
\end{solution}

\begin{solution}{exo:g12:climate-change:14}
$\dfrac{\dd F}{\dd C} = \dfrac{F_{2\times}}{C \ln 2}$. عند \qty{280}{ppm}:
$3.93/(280 \times 0.6931) \approx \qty{0.0202}{W/m^2}$ لكل جزء من المليون؛ وعند
\qty{420}{ppm}: $\approx \qty{0.0135}{W/m^2}$ لكل جزء من المليون. فكل جزء إضافي من المليون يضيف
\omterm{def:g12:climate-change:forcing}{تأثيرًا} أقل كلما ازداد التركيز (أقل بمقدار 1.5 مرة عند 420 منه عند 280)،
لكن التركيز يظل يرتفع، فيظل \omterm{def:g12:climate-change:forcing}{التأثير} يزداد.
\end{solution}

\begin{solution}{exo:g12:climate-change:15}
$3.93/3.22 \approx \qty{1.2}{\degreeCelsius}$ من دون \omterm{def:g12:climate-change:feedback}{التغذيات الراجعة}؛
و$3.93/1.16 \approx \qty{3.4}{\degreeCelsius}$ معها. وعامل
التضخيم: $3.22/1.16 \approx 2.8$.
\end{solution}

\begin{solution}{pb:g12:climate-change:1}
\textbf{1.} $3.37 + 3.51 = \qty{6.88}{ppm}$، أي نحو \qty{7}{ppm}.

\textbf{2.} من الربيع إلى أواخر الصيف، تسحب نباتات نصف الكرة
الشمالي ثاني أكسيد الكربون من الهواء بالتمثيل الضوئي أسرع مما
يعيده التنفس والتحلل: فينخفض التركيز حتى سبتمبر.
وفي الخريف والشتاء ينقلب التوازن، فيعود التركيز إلى الصعود
حتى قمته في مايو.

\textbf{3.} في نصف الكرة الجنوبي يابسة أقل بكثير، ومن ثم غطاء نباتي
أقل بكثير، مما في الشمالي؛ ومعظم سطحه محيط.

\textbf{4.} لأنه عالٍ، يقع فوق طبقة الهواء الملوث قرب
الأرض؛ ولأنه بعيد عن القارات والمدن، فالهواء الذي يأخذ منه العينات ممتزج جيدًا
ويمثل الغلاف الجوي كله، لا مصدرًا محليًا.

\textbf{5.} $(325.68 - 316.91)/10 \approx \qty{0.88}{ppm}$ في السنة.

\textbf{6.} $(427.35 - 401.01)/10 \approx \qty{2.63}{ppm}$ في السنة.

\textbf{7.} $2.634/0.877 \approx 3.0$: لقد تضاعف معدل الزيادة ثلاث مرات.

\textbf{8.} $(427.35 - 278)/278 \approx 0.537$: أي نحو 54\,\%.

\textbf{9.} $427.35 - 300 \approx \qty{127}{ppm}$.

\textbf{10.} $F = 3.93 \times \ln(427.35/278)/\ln 2 \approx
\qty{2.44}{W/m^2}$.

\textbf{11.} $F(316.91) = 3.93 \times \ln(316.91/278)/\ln 2 \approx
\qty{0.74}{W/m^2}$؛ والزيادة $2.44 - 0.74 \approx \qty{1.70}{W/m^2}$.

\textbf{12.} كل \unit{W/m^2} إضافي يضرب التركيز في
$2^{1/3.93} \approx 1.193$: $427.35 \times 1.193 \approx \qty{510}{ppm}$.

\textbf{13.} $(509.8 - 427.35)/2.634 \approx 31$ سنة.

\textbf{14.} $2 \times 278 = \qty{556}{ppm}$؛ \omterm{def:g12:climate-change:forcing}{والتأثير} $F_{2\times} =
\qty{3.93}{W/m^2}$.

\textbf{15.} بتعريف الحساسية، \qty{3}{\degreeCelsius}
(على نحو مرجح من 2.5 إلى \qty{4}{\degreeCelsius}).

\textbf{16.} $\lambda = 3.93/3 \approx \qty{1.31}{W.m^{-2}.\degreeCelsius^{-1}}$.
ومجموع \omterm{def:g12:climate-change:feedback}{التغذيات الراجعة}، \qty{1.16}{W.m^{-2}.\degreeCelsius^{-1}}، يعطي
$3.93/1.16 \approx \qty{3.4}{\degreeCelsius}$، أي أكثر قليلًا؛ وأفضل تقدير
في التقييم يجمع هذه الفيزياء مع الاحترار المرصود والمناخات
الماضية، والقيمتان كلتاهما تقعان في المدى المرجح.

\textbf{17.} يخزن المحيط معظم الطاقة الإضافية ويسخن ببطء: فلا يبلغ
السطح احتراره عند التوازن إلا بعد عقود إلى قرون.

\textbf{18.} لأنه يفترض رواية واحدة للانبعاثات المقبلة (إبقاء زيادة
السنوات العشر الأخيرة إلى الأبد): وهذا \omterm{def:g12:climate-change:scenario}{سيناريو}. أما الإسقاط فيحسب
استجابة المناخ ودورة الكربون \omterm{def:g12:climate-change:scenario}{لسيناريو} انبعاثات
بنماذج فيزيائية، ويعطي مدى وأرجحيات.

\textbf{19.} $(556 - 427.35)/2.634 \approx 48.8$ سنة بعد 2025: أي نحو
عام \textbf{2074}.
\end{solution}
